\begin{table}[htbp]
\centering
\caption{Robustness: Leave-One-Out and Alternative Specifications}
\label{tab:robustness}
\begin{tabular}{lcc}
\hline\hline
 & Log Reg. & Govt.\ Eff. \\
\hline
\textit{Panel A: Baseline} & & \\
\quad Main estimate & 0.197 & 0.374 \\
\quad & (0.145) & (0.092) \\
 & & \\
\textit{Panel B: Leave-one-out (log reg.)} & & \\
\quad Range & [0.132, 0.307] & --- \\
 & & \\
\textit{Panel C: Alternative specifications} & & \\
\quad Restricted donors (Caucasus + C.\ Asia) & -0.141 & --- \\
\quad & (0.209) & \\
\quad Placebo treatment (2010) & -0.322 & --- \\
\quad & (0.085) & \\
\quad GDP per capita (placebo outcome) & \multicolumn{2}{c}{-2,753} \\
\quad & \multicolumn{2}{c}{(526)} \\
 & & \\
\textit{Panel D: Permutation inference} & & \\
\quad Rank of Azerbaijan & 8/10 & --- \\
\quad Permutation $p$-value & 0.800 & --- \\
\hline\hline
\end{tabular}
\par\vspace{0.3em}
{\footnotesize \textit{Notes:} Panel A reproduces baseline estimates from Table~\ref{tab:did_results}. Panel B drops one donor country at a time. Panel C tests alternative specifications: restricting to 4 Caucasus and Central Asian donors; assigning placebo treatment to 2010 (pre-ASAN); using GDP per capita as a placebo outcome (the negative coefficient reflects Azerbaijan's 2014--2016 oil price shock). Panel D ranks Azerbaijan's ATT against placebo ATTs from each donor country. Standard errors clustered at the country level.}
\end{table}
